% Chapitre 5 — Conception et architecture
\chapter{Conception et architecture du système}
\label{ch:conception}

\section{Architecture logique}

L'architecture suit un modèle \textbf{3-tiers} enrichi d'une couche IA :

\begin{enumerate}[leftmargin=*]
  \item \textbf{Présentation} : application React (Explain + Chat) ;
  \item \textbf{Application} : API FastAPI + orchestrateur LangGraph ;
  \item \textbf{Domaine ML} : modèles classiques, séquentiels, graphe, score KYC ;
  \item \textbf{Données} : CSV bruts, Parquet, features, index RAG, artefacts Joblib/PyTorch.
\end{enumerate}

\todofig{Diagramme de déploiement — Client React, API Uvicorn, Ollama, fichiers locaux}

\section{Structure du dépôt}

\begin{lstlisting}[caption={Arborescence principale du projet},label={lst:tree}]
Talys_pfe/
|-- data/raw/           # CSV (clients, credits, ...)
|-- data/processed/     # Parquet nettoyes
|-- data/features/      # features.parquet
|-- src/
|   |-- api/            # FastAPI (main.py, schemas.py)
|   |-- agent/          # LangGraph orchestrator
|   |-- rag/            # Index TF-IDF + retrieval
|   |-- kyc/            # Score KYC
|   |-- models/         # train, sequential, graph
|   |-- features/       # engineering.py
|   |-- data/           # loader, cleaner
|   |-- llm/            # client Ollama
|-- frontend/           # React + Vite
|-- models/             # Artefacts (.joblib, .pt)
|-- reports/            # Figures, knowledge_base, generated
\end{lstlisting}

\section{Workflow LangGraph}

L'orchestrateur \texttt{CreditAgentOrchestrator} implémente le graphe d'états suivant :

\begin{lstlisting}[language={},caption={Workflow LangGraph (PlantUML activité)},label={lst:langgraph}]
START -> ExtractCIN -> DetectIntent -> ChoosePath
  |-> RunClassic ----\
  |-> RunSequential -+-> RetrieveDocs -> WriteReport -> Respond -> END
  |-> RunGraph ------/
  |-> RunFullReport (3 modeles)
  |-> RunCompareModels (3 modeles)
\end{lstlisting}

\subsection{Nœuds et responsabilités (SRP)}

\begin{table}[H]
  \centering
  \caption{Nœuds LangGraph}
  \begin{tabular}{|l|p{9cm}|}
    \hline
    \textbf{Nœud} & \textbf{Responsabilité} \\
    \hline
    ExtractCIN & Regex CIN + credit\_id optionnel \\
    \hline
    DetectIntent & Mots-clés FR/EN (séquentiel, graphe, rapport, comparer) \\
    \hline
    ChoosePath & Mapping intent $\rightarrow$ route \\
    \hline
    RunClassic/Sequential/Graph & Appel endpoints internes (LLM désactivé) \\
    \hline
    RetrieveDocs & Multi-query RAG (TF-IDF) \\
    \hline
    WriteReport & Prompt LLM + fallback Markdown \\
    \hline
    Respond & Assemblage réponse finale \\
    \hline
  \end{tabular}
\end{table}

\section{Diagramme de classes (vue logique)}

\begin{lstlisting}[language={},caption={Diagramme de classes simplifié (PlantUML)},label={lst:class}]
@startuml
package Frontend {
  class ReactApp
  class ApiClient
}
package Backend {
  class FastAPIApp
}
package Orchestration {
  class CreditAgentOrchestrator
  class AgentState
}
package Scoring {
  class ClassicScoring
  class SequentialScoring
  class GraphScoring
}
package RAG {
  class RagIndex
  class RagRetriever
}
ReactApp --> FastAPIApp
FastAPIApp --> CreditAgentOrchestrator
CreditAgentOrchestrator --> ClassicScoring
CreditAgentOrchestrator --> SequentialScoring
CreditAgentOrchestrator --> GraphScoring
CreditAgentOrchestrator --> RagRetriever
@enduml
\end{lstlisting}

\section{Endpoints API principaux}

\begin{longtable}{|l|l|p{6cm}|}
  \caption{Endpoints FastAPI exposés (Swagger)} \\
  \hline
  \textbf{Méthode} & \textbf{Route} & \textbf{Description} \\
  \hline
  \endfirsthead
  \hline
  \textbf{Méthode} & \textbf{Route} & \textbf{Description} \\
  \hline
  \endhead
  GET & /health & Santé API + nom modèle \\
  \hline
  POST & /explain/by-cin & Scoring classique + LLM \\
  \hline
  POST & /explain/sequential/by-cin & Scoring séquentiel multi-crédits \\
  \hline
  POST & /explain/graph/by-cin & Scoring GraphSAGE \\
  \hline
  POST & /chat & Agent LangGraph \\
  \hline
  POST & /report/by-cin & Rapport Markdown + sources RAG \\
  \hline
  POST & /report/by-cin/download & Export MD ou PDF \\
  \hline
\end{longtable}

Les routes \api{/predict*} existent mais sont masquées (\texttt{include\_in\_schema=False}) car l'usage métier passe par \api{/explain/*} et \api{/chat}.

\section{Choix de conception clés}

\subsection{Séparation scoring / explication}

Les scores numériques sont \textbf{toujours} produits par les modèles ML déterministes. Le LLM et le RAG interviennent \textbf{uniquement} en aval pour rédiger des textes. Cette séparation garantit l'auditabilité réglementaire.

\subsection{Entrée CIN seule}

Les endpoints \api{/explain/*/by-cin} chargent automatiquement les features depuis les tables lookup (\texttt{clients.csv}, \texttt{credits.csv}, \texttt{features.parquet}), simplifiant l'intégration ICM future.

\subsection{Index RAG curé}

Seuls les documents métier (\texttt{knowledge\_base/}, \texttt{criteres\_score\_et\_crispdm.md}) sont indexés, excluant les rapports générés et fichiers de debug.

\section{Sécurité et conformité (prototype)}

\begin{itemize}[leftmargin=*]
  \item Pas de données PII réelles en démonstration ;
  \item CORS activé pour le frontend local ;
  \item Disclaimer dans les rapports : décision humaine obligatoire ;
  \item Pas de conseil juridique généré automatiquement.
\end{itemize}

Perspectives : authentification JWT, chiffrement des logs, hébergement sécurisé.
